% ============================================================
\chapter{Fixed Point Theory in Advanced Topology}
\label{ch:fixed_points}
% ============================================================

\begin{intuition}
Fixed point theorems play a central role in theoretical computer science and domain theory.
Unlike the classical theorems of Banach or Brouwer, which rely on metric or classical
topological structures, the theorems presented here exploit order and directed completeness.
Tarski's theorem guarantees the existence of fixed points in complete lattices, Kleene's
theorem constructs them effectively by iteration, and Pataraia's theorem extends these
results to dcpos without any countability hypothesis.
\end{intuition}

% ------------------------------------------------------------
\section{Tarski's Theorem for Complete Lattices}
% ------------------------------------------------------------

\begin{definition}[Complete lattice]
\index{complete lattice}
A partially ordered set $(L, \leq)$ is a \emph{complete lattice} if every subset
$S \subseteq L$ has a supremum $\bigvee S$ and an infimum $\bigwedge S$. In particular,
$L$ has a least element $\bot = \bigvee \emptyset$ and a greatest element
$\top = \bigwedge \emptyset$.
\end{definition}

\begin{theorem}[Knaster--Tarski]
\label{thm:knaster-tarski}
\index{Knaster--Tarski theorem}
Let $(L, \leq)$ be a complete lattice and $f\colon L \to L$ a monotone function. Then:
\begin{enumerate}
  \item $f$ has a least fixed point
    $\mathrm{lfp}(f) = \bigwedge \{x \in L : f(x) \leq x\}$.
  \item $f$ has a greatest fixed point
    $\mathrm{gfp}(f) = \bigvee \{x \in L : x \leq f(x)\}$.
  \item The set of fixed points $\mathrm{Fix}(f)$ is a complete lattice.
\end{enumerate}
\end{theorem}

\begin{proof}
(1) Let $P = \{x \in L : f(x) \leq x\}$ and $a = \bigwedge P$. For every $x \in P$,
$a \leq x$, so $f(a) \leq f(x) \leq x$ by monotonicity. Thus $f(a)$ is a lower bound
of $P$, hence $f(a) \leq a$. This shows $a \in P$.

By monotonicity, $f(f(a)) \leq f(a)$, so $f(a) \in P$, hence $a \leq f(a)$. Combined
with $f(a) \leq a$, we get $f(a) = a$.

If $b$ is another fixed point, then $b \in P$ so $a \leq b$: $a$ is indeed the least
fixed point.

(2) Dual of (1).

(3) Let $S \subseteq \mathrm{Fix}(f)$. The set
$\{x \in L : x \geq s \; \forall s \in S, \; f(x) \leq x\}$ is nonempty (it contains
$\top$). Its infimum is a fixed point of $f$ that is the supremum of $S$ in
$\mathrm{Fix}(f)$.
\end{proof}

\begin{remark}
The Knaster--Tarski theorem requires no continuity hypothesis: monotonicity alone suffices.
However, it does not provide a constructive method for computing the fixed point.
\end{remark}

\begin{example}[Application to language equations]
Let $\Sigma$ be an alphabet and $\mathcal{P}(\Sigma^*)$ the lattice of languages over
$\Sigma$. A system of language equations $X_i = f_i(X_1, \ldots, X_n)$ where the $f_i$
are monotone has a least solution, by Tarski's theorem applied to the lattice
$\mathcal{P}(\Sigma^*)^n$.
\end{example}

% ------------------------------------------------------------
\section{Kleene's Fixed Point Theorem}
% ------------------------------------------------------------

\begin{definition}[Pointed dcpo]
\index{pointed dcpo}
A \emph{pointed dcpo} is a dcpo $(D, \leq)$ that has a least element $\bot$.
\end{definition}

\begin{theorem}[Kleene]
\label{thm:kleene}
\index{Kleene's theorem}
Let $(D, \leq)$ be a pointed dcpo and $f\colon D \to D$ a Scott-continuous function. Then
$f$ has a least fixed point, given by:
\[
  \mathrm{lfp}(f) = \bigvee_{n \geq 0} f^n(\bot).
\]
\end{theorem}

\begin{proof}
First we show that the sequence $(f^n(\bot))_{n \geq 0}$ is increasing. We have
$\bot \leq f(\bot)$ since $\bot$ is the least element. By induction, if
$f^n(\bot) \leq f^{n+1}(\bot)$, then by monotonicity,
$f^{n+1}(\bot) \leq f^{n+2}(\bot)$.

The set $\{f^n(\bot) : n \geq 0\}$ is directed (it is a chain), so its supremum
$a = \bigvee_n f^n(\bot)$ exists in $D$.

By Scott-continuity:
\[
  f(a) = f\!\left(\bigvee_{n \geq 0} f^n(\bot)\right) = \bigvee_{n \geq 0} f^{n+1}(\bot) = a.
\]
Hence $a$ is a fixed point.

If $b$ is another fixed point, then $\bot \leq b$ and by induction $f^n(\bot) \leq b$
for all $n$, whence $a = \bigvee_n f^n(\bot) \leq b$.
\end{proof}

\begin{attention}
Kleene's theorem requires Scott-continuity (not merely monotonicity). A monotone function
on a pointed dcpo may not have a least fixed point obtained by iteration from $\bot$.
\end{attention}

\begin{corollary}[Transfinite iteration]
If $f\colon L \to L$ is a monotone function on a complete lattice, the least fixed point
is reached by transfinite iteration:
\[
  f^\alpha(\bot) = \begin{cases}
    f(f^{\alpha-1}(\bot)) & \text{if } \alpha \text{ is a successor ordinal}, \\
    \bigvee_{\beta < \alpha} f^\beta(\bot) & \text{if } \alpha \text{ is a limit ordinal}.
  \end{cases}
\]
There exists an ordinal $\alpha_0$ such that $f^{\alpha_0}(\bot) = \mathrm{lfp}(f)$.
\end{corollary}

\begin{example}[Semantics of loops]
The semantics of a \texttt{while b do c} loop is the least fixed point of the operator
$\Phi\colon [S \to S_\bot] \to [S \to S_\bot]$ defined by:
\[
  \Phi(g)(s) = \begin{cases}
    g(\llbracket c \rrbracket(s)) & \text{if } \llbracket b \rrbracket(s) = \mathtt{true}, \\
    s & \text{if } \llbracket b \rrbracket(s) = \mathtt{false}, \\
    \bot & \text{otherwise}.
  \end{cases}
\]
By Kleene's theorem, $\mathrm{lfp}(\Phi) = \bigvee_n \Phi^n(\bot_{S \to S_\bot})$.
\end{example}

% ------------------------------------------------------------
\section{Pataraia's Theorem}
% ------------------------------------------------------------

\begin{theorem}[Pataraia]
\label{thm:pataraia}
\index{Pataraia's theorem}
Let $(D, \leq)$ be a pointed dcpo and $f\colon D \to D$ a monotone function (not
necessarily continuous) such that $f(x) \geq x$ for all $x$ (\emph{inflationary}). Then
$f$ has a fixed point.
\end{theorem}

\begin{proof}
Define the collection $\mathcal{I}$ of all subsets $S$ of $D$ such that:
\begin{enumerate}
  \item $\bot \in S$;
  \item $S$ is closed under $f$;
  \item $S$ is closed under directed suprema (in $D$).
\end{enumerate}
Let $C = \bigcap \mathcal{I}$ be the smallest such set. One shows that $C$ is a chain (by
verifying that the set of elements comparable to every element of $C$ satisfies the three
conditions). The supremum $c = \bigvee C$ exists in $D$ and satisfies $f(c) \leq c$ (since
$f(c) \in C$ by closure, whence $f(c) \leq c$). Since $f$ is inflationary, $c \leq f(c)$,
hence $f(c) = c$.
\end{proof}

\begin{remark}
Pataraia's theorem generalizes the Bourbaki--Witt theorem (which assumes that every chain
has a supremum). The advantage of Pataraia's theorem is that it requires only directed
completeness, not chain completeness.
\end{remark}

\begin{corollary}
Every pointed dcpo in which every monotone inflationary function has a fixed point is a
complete lattice.
\end{corollary}

% ------------------------------------------------------------
\section{Least Fixed Points in Domain Theory}
% ------------------------------------------------------------

\begin{theorem}[Existence of the least fixed point]
\label{thm:lfp-domain}
Let $D$ be a pointed continuous domain and $f\colon D \to D$ a Scott-continuous function.
Then:
\begin{enumerate}
  \item The least fixed point $\mathrm{lfp}(f)$ exists and is given by Kleene's formula
    $\mathrm{lfp}(f) = \bigvee_n f^n(\bot)$.
  \item If moreover $D$ is an effectively presented algebraic domain and $f$ is
    effectively given, then $\mathrm{lfp}(f)$ is computable.
  \item The operator $\mathrm{lfp}\colon [D \to D] \to D$ is itself Scott-continuous.
\end{enumerate}
\end{theorem}

\begin{proof}
(1) This is Kleene's theorem (Theorem~\ref{thm:kleene}).

(2) The finite approximations $f^n(\bot)$ are computable, and the compact elements below
$\mathrm{lfp}(f)$ are enumerable.

(3) Let $(f_i)_{i \in I}$ be a directed family in $[D \to D]$ and $g = \bigvee_i f_i$.
Then:
\[
  \mathrm{lfp}(g) = \bigvee_n g^n(\bot) = \bigvee_n \left(\bigvee_i f_i\right)^n(\bot)
  = \bigvee_i \bigvee_n f_i^n(\bot) = \bigvee_i \mathrm{lfp}(f_i).
\]
The interchange of suprema is justified by the Scott-continuity of composition and
application.
\end{proof}

% ------------------------------------------------------------
\section{Applications to Recursive Definitions}
% ------------------------------------------------------------

\begin{proposition}[Recursive function definitions]
Let $D$ be a pointed dcpo. A recursive definition of the form
\[
  f(x) = \Phi(f, x)
\]
where $\Phi\colon [D \to D] \times D \to D$ is Scott-continuous, has a least solution
$f^* \in [D \to D]$, given by
$f^* = \mathrm{lfp}(\lambda g.\, \lambda x.\, \Phi(g,x))$.
\end{proposition}

\begin{example}[Factorial]
The recursive definition
\[
  \mathtt{fact}(n) = \begin{cases} 1 & \text{if } n = 0, \\ n \cdot \mathtt{fact}(n-1) & \text{otherwise} \end{cases}
\]
corresponds to $\mathtt{fact} = \mathrm{lfp}(\Phi)$ where
$\Phi(g)(n) = \text{if } n = 0 \text{ then } 1 \text{ else } n \cdot g(n-1)$. The
operator $\Phi$ is Scott-continuous on $[\N_\bot \to \N_\bot]$.
\end{example}

\begin{theorem}[Recursive domain equations]
\index{domain equations}
If $F$ is a locally continuous functor on the category of continuous domains with
retraction pairs, then the domain equation
\[
  D \cong F(D)
\]
has a solution, obtained as a projective limit (and also inductive colimit) of the
sequence $D_0 \leftarrow D_1 \leftarrow \cdots$ where $D_{n+1} = F(D_n)$.
\end{theorem}

\begin{keyformulas}[title=\textbf{Summary of fixed point theorems}]
\begin{center}
\renewcommand{\arraystretch}{1.3}
\begin{tabular}{|l|l|l|l|}
\hline
\textbf{Theorem} & \textbf{Structure} & \textbf{Hypothesis on $f$} & \textbf{Result} \\
\hline
Knaster--Tarski & Complete lattice & Monotone & $\mathrm{Fix}(f)$ complete lattice \\
Kleene & Pointed dcpo & Scott-continuous & $\mathrm{lfp}(f) = \bigvee f^n(\bot)$ \\
Pataraia & Pointed dcpo & Monotone inflat. & Existence of fixed point \\
Bourbaki--Witt & Chain-complete & Inflationary & Existence of fixed point \\
\hline
\end{tabular}
\end{center}
\end{keyformulas}

% ------------------------------------------------------------
\section{Exercises}
% ------------------------------------------------------------

\begin{exercise}
Let $L = \mathcal{P}(\N)$ be the lattice of subsets of $\N$. Determine the least and
greatest fixed points of $f(A) = \N \setminus A$.
\end{exercise}

\begin{exercise}
Show that in Kleene's theorem, the sequence $(f^n(\bot))$ can stabilize after finitely
many steps if and only if the least fixed point is a compact element of the dcpo.
\end{exercise}

\begin{exercise}
Give an example of a pointed dcpo $D$ and a monotone function $f\colon D \to D$ such that
the sequence $(f^n(\bot))$ does not converge to the least fixed point of $f$ (i.e., such
that $\bigvee_n f^n(\bot) < \mathrm{lfp}(f)$).
\end{exercise}

\begin{exercise}
Let $D$ be a pointed dcpo and $f, g\colon D \to D$ two Scott-continuous functions with
$f \leq g$ pointwise. Show that $\mathrm{lfp}(f) \leq \mathrm{lfp}(g)$.
\end{exercise}

\begin{exercise}[Bekic's theorem]
Let $D, E$ be pointed dcpos and $f\colon D \times E \to D$, $g\colon D \times E \to E$
Scott-continuous functions. Show that the least fixed point of
$(f,g)\colon D \times E \to D \times E$ is given by
$(\mathrm{lfp}(\hat{f}), \mathrm{lfp}(\hat{g}))$ where $\hat{f}$ and $\hat{g}$ are
defined by mutual substitution.
\end{exercise}

\begin{exercise}
Let $\Sigma = \{a, b\}$ and $\mathcal{P}(\Sigma^*)$ the lattice of languages. The
operator $\Phi(X) = \{\varepsilon\} \cup \{a\} \cdot X \cdot \{b\}$ is monotone. Compute
$\mathrm{lfp}(\Phi)$ using Kleene's theorem.
\end{exercise}

\begin{exercise}
Show that Pataraia's theorem implies Zorn's lemma: every inductive partially ordered set
(every chain is bounded) has a maximal element.
\emph{Hint: consider a pointed dcpo and a choice function.}
\end{exercise}
